% =====================================================================
% intro.tex — DRAFTED at rb-049 (RB-049), 2026-05-30.
%
% Third bookend; companion of sections/abstract.tex (rb-047, RB-047)
% and sections/limitations.tex (rb-048, RB-048).  Mirrors the abstract's
% voice (distributional / graded / conditional by default) and sets up
% the §3.3 unifying-reframe contrast that the §limitations §7.6 'what
% the rebuild does not claim' paragraph closes on.
%
% Strength check (mission §3, CLAIM_LEDGER.md): every quantitative
% claim below is at the licensed band / distributional / conditional
% strength already wired in the body sections.  Original-paper
% categorical phrasings (CF floor [0.60, 0.96]; "negligible regardless
% of other parameters"; uniform no-inversion; A1 upper bound on VDA)
% appear here only as the framing being retracted, not stated.
% =====================================================================

\section{Introduction}
\label{sec:intro}

In a cued change-detection task whose cue carries both \emph{value}
$\val \ge 1$ and \emph{validity} $\valid \in [1/\Nloc, 1]$, an
observer adapting to the value contingency has more than one lever
available.  They may shift decision criteria so that responses at the
high-value location are biased toward report; they may re-allocate
spatial attention so that the perceptual evidence at the high-value
location is itself more reliable; and, when noise is correlated
across locations, they may decorrelate the encoded population code so
that the joint evidence drawn from the unattended array is less
redundant.  The original \emph{value-directed attention} (VDA) model
asks a sharp normative question on top of these levers: how much of
the behavioural adaptation to value is achieved by re-allocating
attention toward the high-value location versus by shifting decision
criteria, and in what parameter regime does attention re-allocation
become the dominant mechanism?  This paper preserves that question
verbatim in spirit.  What changes is the mathematical description used
to answer it.

The inherited model couples value-driven sensitivity changes to a
single attentional asymmetry parameter and decomposes the expected
reward across nested policies that toggle each lever on or off.  Its
quantitative claims, however, tend to overstate what their underlying
simulations license: a central or peak value is repeatedly reported
as a uniform floor, a directional inequality, or a
``regardless of other parameters'' constant.  Re-evaluating those
claims against the model's own definitions reveals a recurring
mismatch.  The reported criterion-fraction interval $[0.60, 0.96]$ is
in fact $[0.30, 1.00]$ across the same $4{,}410$-cell sweep with
median $0.76$ and a $\sim\!13\%$ tail below $0.60$; the
``high-validity paradigms show negligible VDA regardless of other
parameters'' design recommendation holds at $\valid \ge 0.95$ but
not at $\valid \approx 0.80$ once cross-location correlation is
admitted; the ``no inversion'' result is true everywhere in the cued
regime but not in the anti-cue regime $\valid < 1/\Nloc$; and the
self-characterisation of independence as ``an upper bound on VDA'' is
sign-ambiguous --- correlation \emph{suppresses} VDA in the
cost-dominant regime and \emph{amplifies} it in the
benefit-dominant tail.  None of these reframings overturns the
underlying mechanism the original paper identified.  Each one
restates a central tendency, a graded boundary, or a conditional
theorem in place of an unconditional one.  Reporting distributions,
quantiles, contour bands, and explicit conditions \emph{instead of}
floors, ``never,'' and ``regardless of'' is the corrective this paper
applies uniformly: distributional, graded, or conditional by default.

Restating the model in this voice also exposes a missing lever.  The
inherited two-channel decomposition treats criterion shifts and
sensitivity re-allocation as the only adaptive mechanisms; the
empirically dominant decorrelation channel
\cite{CohenMaunsell2009,
RuffCohen2016,
Srinath2021} is excluded by an
independence assumption that the model bundles into its
no-false-alarm probability.  Section~\ref{sec:model} promotes
cross-location correlation $\corr$ to a first-class model parameter
through an exact one-factor Gauss--Hermite reduction of the
equicorrelated-Gaussian orthant integral, recovering the inherited
independent-noise model in the limit $\corr \to 0$ to floating-point
identity (Equation~\eqref{eq:pnofa-rho},
Equation~\eqref{eq:rho-zero-recovery}).  This yields the rebuilt
paper's central reframe: three levers, not two
(Definition~\ref{def:three-levers}; criterion shift,
sensitivity re-allocation, decorrelation).  The inherited model's
self-characterisation of independence as an upper bound on VDA is
replaced by a more careful statement: what independence actually
upper-bounds is the criterion fraction $\CF$ ---
$\CF(\corr) \le \CF(0)$ pointwise in variant~A by a Slepian
monotonicity argument (Appendix~\ref{sec:appendix-deriv-a1}) --- while
$\partial \VDA / \partial \corr$ flips sign near $\Rsens \approx 0.5$
at the headline cell and sweeps a cell-wise crossover up to
$\Rsens \approx 0.79$ across the full $4{,}410$-cell sweep
(Section~\ref{sec:model-upper-bound},
Table~\ref{tab:a1cw-summary}).

Under the corrective voice, the four headline C-row results recover
in distributional, graded, and conditional forms.  The criterion
fraction is concentrated in $[0.30, 1.00]$ with median $0.76$
(Section~\ref{sec:results-c1}, Table~\ref{tab:cf-distribution});
$\VDA$ is non-monotonic in $\Rsens$ with a closed-form lower edge
$\rdagger(\val) = K_u(\val) / [(\Nloc - 1) K_c(\val)]$
(Equation~\eqref{eq:r-dagger}) that extends to $\corr > 0$ through
$\corr$-aware gradient quadratures
(Section~\ref{sec:appendix-deriv-c2},
Proposition~\ref{prop:r-dagger-rho}) and is invariant under the
benefit/cost conservation order
(Proposition~\ref{prop:r-dagger-invariance}); the ``narrow regime''
becomes a quantitative iso-$\VDA$ contour band
(Section~\ref{sec:results-c3}, Table~\ref{tab:c3-highV-probe}) with
hierarchical $\valid$ thresholds ($\valid \ge 0.95$ at the grid floor
under any $(\Rsens, \corr)$; $\valid \ge 0.80$ only at $\corr = 0$;
$\valid \ge 0.60$ fails); and the no-inversion result becomes a
conditional theorem with the explicit boundary
$\valid \ge 1 / [(\Nloc - 1) \val + 1]$
(Section~\ref{sec:results-c4}, Equation~\eqref{eq:value-weight})
that, when crossed, generates a new falsifiable prediction --- the
\emph{anti-cue inversion}: under counter-predictive cues
($\valid < 1/\Nloc$) at $\Nloc = 4$, the normatively optimal
allocation places \emph{less} attention at the high-value location
in $36.1\%$ of probed cells (Table~\ref{tab:c4-anticue}).  Symmetric
recovery at $\Rsens = 1$ is reaffirmed as a universal real-number
identity (Appendix~\ref{sec:appendix-c5},
Proposition~\ref{prop:c5-realnumber}), conservation-form-invariant by
construction (Corollary~\ref{cor:a3d-c5-invariance}).

Three further mechanisms the inherited model held fixed at
homogeneous values return as explicit extensions.
Within-display heterogeneity of $\Rsens$ is a bounded perturbation
that leaves the $C_2$ peak coordinates and the contested-CF corner
essentially invariant (Section~\ref{sec:extensions-a2},
Table~\ref{tab:a2-rb021-summary}).  The benefit/cost conservation
rule $\benefit + \cost = 2$ is one point in the power-mean
conservation family
(Equation~\eqref{eq:conservation-family},
Equation~\eqref{eq:beta-gamma-of-p}); the rebuilt paper reports
every headline number as a band over the family at
$p \in \{0, 0.5, 1\}$ with the per-cell monotonicity
$\Delta \CF(\corr = 0) \le 0$ stated as
Theorem~\ref{thm:delta-cf-monotone}
(Section~\ref{sec:extensions-a3}).  The $N$-dimensional uncued
allocation policy is innocuous at the inherited additive optimum but
exhibits a new conditional binding in a benefit-dominant
symmetric-stress corner under multiplicative conservation, with a
$32\%$ $\corr$-amplification at the focal cell
(Section~\ref{sec:extensions-a8}).  Per-location decision-noise
coupling --- a candidate fourth lever
--- is named and held at the live status of its underlying claim
(Section~\ref{sec:limitations-a6}) rather than wired into the model.

The remainder of the manuscript follows this layout.
Section~\ref{sec:model} states the rebuilt model with the
three-lever decomposition and the $\corr$-channel reduction.
Section~\ref{sec:results} reports the C-row headline results in the
distributional / graded / conditional voice, with each subsection
naming the inherited categorical claim it replaces.
Section~\ref{sec:extensions} runs the three lever extensions
(A2, A3, A8) at empirical band strength.
Section~\ref{sec:limitations} catalogues the explicit scope of
every claim --- in particular, what the rebuilt model does
\emph{not} claim about neural implementation,
empirical-$\corr$ prediction, observer-deviation, or attention
dynamics.  The appendix collects the formal derivations
(A1 $\corr$ channel, $C_2$ closed forms,
$C_5$ symmetric recovery, A3 power-mean conservation).
Every quantitative claim in the body is backed by a simulation in
\texttt{Rebuild/sims/} with a recovered limit, a deterministic
output hash, and a captioned figure; every novel proposition is
derived in \texttt{Rebuild/derivations/}.  Reproducibility is
end-to-end byte-identical, and nothing in this paper is stated more
strongly than its row of \texttt{CLAIM\_LEDGER.md} licenses.
